% =============================================================================
% Appendix C — Open Problems and Research Directions
% =============================================================================
\section{Open Problems and Research Directions}\label{app:directions}

Every \textbf{Open} line in this book is a claim that something has not been
done. This appendix turns those claims into a work programme, ordered by what we
judge to be the ratio of value to effort, and written so that a reader could
start on any one of them tomorrow.

Each direction states four things, and the fourth is the one usually omitted:

\begin{itemize}[leftmargin=*,itemsep=1pt,topsep=2pt]
  \item \textbf{The question}, in a form that admits an answer.
  \item \textbf{Where it comes from} --- the ledger line that records the gap.
  \item \textbf{First step}, concrete enough to begin.
  \item \textbf{What would refute it}, or in the case of a measurement, what
        outcome would make the direction not worth pursuing. A research
        programme with no such outcome is a plan to generate confirmations.
\end{itemize}

A note on honesty about priority. These are ranked by our judgement and nothing
more; we have not surveyed the field to check that items~1--3 are unattempted,
and a reader should assume some are already in progress elsewhere.

\subsection*{D1. Measure the synthesis ratio}

\textbf{The question.} What fraction of the genome does \texttt{pol2-P301R}
copy, relative to the exonuclease-null enzyme, \emph{in vivo}?

\textbf{Where it comes from.} Section~\ref{sec:mmr}'s ledger and
Remark~\ref{rem:transport}. This book's replacement mechanism --- that repair
capacity must fall --- rests on the hypothesis
$p_{\mathrm{mut}} \le p_{\mathrm{null}}$, which is a statement about error load
delivered per replicated genome. What is measured is error rate per nucleotide
incorporated. \lean{transport_needs_synthesis_bound} names the missing factor
exactly: the ratio $s_{\mathrm{mut}}/s_{\mathrm{null}}$.

\textbf{Why it is first.} Exercise~11.5 makes the stakes numerical. The in vitro
accuracy advantage is cancelled once $s_{\mathrm{mut}}/s_{\mathrm{null}} > 155/59
\approx 2.63$, and the entire in vivo ratio would be accounted for without any
repair effect at $\approx 137$. The first figure is not absurd for a hyperactive
polymerase; the second is. So the measurement discriminates sharply between
``the capacity story is right'', ``it is partly right'', and ``it is unnecessary''
--- and no other single experiment in this book's orbit does that.

\textbf{First step.} Strand-specific measurement of polymerase usage in the two
strains. Ribonucleotide-incorporation mapping (HydEn-seq or Pu-seq in an
\texttt{rnh201}$\Delta$ background) gives the fraction of each strand synthesised
by Pol~$\varepsilon$ against Pol~$\delta$ directly, and both strains are already
constructed and published.

\textbf{What would make this not worth pursuing.} If $s_{\mathrm{mut}} /
s_{\mathrm{null}}$ comes out indistinguishable from $1$, the transport objection
evaporates and Section~\ref{sec:mmr}'s conclusion becomes unconditional ---
which is a good outcome but means the direction is closed in one experiment. The
risk is the opposite: a ratio near $2.6$ would leave the question genuinely
undetermined and require the hypomorphic-MMR experiment of Exercise~5.4 as well.

\subsection*{D2. Ask whether $\Lambda \ge 52$ is physically achievable}

\textbf{The question.} Is there any cooperative mechanism by which proofreading
stages interact strongly enough that removing them is fifty-two-fold more
damaging than their product?

\textbf{Where it comes from.} Section~\ref{sec:nogo}'s ledger, discharged into
\S\ref{sub:interact}. Theorem~\ref{thm:psr3i} prices the one escape route the
no-go theorem leaves, but pricing is not closing: nothing here shows the price
cannot be paid.

\textbf{First step.} This is theory before it is experiment. Replace the single
lumped $\Lambda$ of \eqref{eq:interact} with a stage-resolved interaction ---
the natural candidate is a Hopfield cascade in which the branch rates of stage
$i$ depend on occupancy at stage $i-1$ --- and compute the maximum attainable
$\Lambda$ over physically admissible rate constants. If that maximum has a
closed form, the result is a second no-go theorem, and a stronger one: it would
exclude the interacting class as well.

\textbf{What would refute it.} A construction reaching $\Lambda \ge 52$ with
rates in the measured range for a B-family polymerase. That would reopen the
stage-removal account and would be a far more interesting outcome than
confirming the refutation. Worked Example~11.2 gives the equivalent statement:
each stage would need to be $\approx 3.7$ times leakier in the mutant than its
own measured escape probability, and an enzymologist may well know immediately
whether that is attainable.

\subsection*{D3. A persistence statistic with a reason to work}

\textbf{The question.} Is there a topological statistic of a sequence embedding
that detects biological structure surviving a dinucleotide-preserving null?

\textbf{Where it comes from.} \S\ref{sub:tdanull}. The negative result there is
specific: $H_1$ total persistence of a 3D chaos-game cloud, on three coding
sequences, against three nulls. Worked Example~11.1 is explicit that this does
not generalise to ``topology does not work on sequence''.

\textbf{Why the current answer is probably no, and what would change it.} The
chaos game is a contraction onto a self-similar attractor, so its topology is
dominated by the attractor and only perturbed by the sequence. Any statistic
computed on that cloud is fighting the embedding. The promising moves change the
embedding rather than the statistic:
\begin{itemize}[leftmargin=*,itemsep=1pt,topsep=2pt]
  \item \textbf{Embed a signal, not the letters.} Replication timing, GC skew or
        nucleosome occupancy along the chromosome is a genuine spatial series;
        sublevel-set persistence of such a track has real features (origins,
        domain boundaries) that a circular-permutation null can test.
  \item \textbf{Embed a contact map.} Hi-C data is already a metric structure,
        and persistent homology of its Rips filtration asks a question ---
        how many independent loops at what scale --- that no simple statistic
        answers.
  \item \textbf{Compare rather than measure.} Persistence landscapes of
        orthologous sequences, where the embedding artefact is common to both and
        cancels in the difference. \S\ref{sub:tdanull} runs three genes but
        compares each only to its own shuffles, never to each other.
\end{itemize}

\textbf{What would refute it.} A circular-permutation null that the candidate
statistic fails to beat on a track where the features are independently known
--- origins from Repli-Seq, say. That is the honest test, and it should be run
before any claim, not after.

\subsection*{D4. The hypomorphic mismatch-repair experiment}

\textbf{The question.} Which of titration, interference, substrate invisibility
or loss of extrinsic Pol~$\delta$ proofreading reduces the repair capacity?

\textbf{Where it comes from.} Section~\ref{sec:mmr}'s ledger and
Exercise~5.4.

\textbf{First step.} Measure \texttt{CAN1} rates for \texttt{pol2-P301R} and
\texttt{pol2-4} across a graded series of MMR activity rather than the null.
Hypomorphic \texttt{msh2} or \texttt{mlh1} alleles, or titratable expression from
a regulated promoter, give $0 < c < c_{\mathrm{null}}$. Under
Assumption~\ref{model:channel}, titration predicts that the capacity consumed
scales with the error load the variant produces, while interference predicts a
reduction independent of load --- so the two give different shapes for escape as
a function of $c$, and a graded series separates them where a null background
cannot.

\textbf{What would make this inconclusive.} If the double mutants are inviable
across the whole graded series, the window closes and the experiment returns
nothing. Bulock \emph{et al.}\ (2020) showing that both MMR and Pol~$\delta$
proofreading are required for viability makes this a real risk, and it should be
checked with a viability titration before the rate measurements are attempted.

\subsection*{D5. Variant effect as a function of repair background}

\textbf{The question.} Does mismatch-repair status modify the relationship
between \texttt{POLE} genotype and tumour mutational burden in patients?

\textbf{Where it comes from.} Section~\ref{sec:clinic}'s ledger and
Exercise~10.4 --- and, more sharply than when that exercise was written, from
the \texttt{V411L} correction of Worked Example~7.1. A variant whose phenotype
appears only when repair is removed is exactly what Section~\ref{sec:mmr}
predicts, and \texttt{V411L} now looks like an instance.

\textbf{Why this is the clinically consequential one.} It is the only direction
here that could change what a clinician does. If MMR status modifies the
variant-to-burden relationship, then interpreting a \texttt{POLE} variant of
uncertain significance requires knowing the tumour's MMR status --- before a
burden has accumulated, which is when the interpretation matters. The base rates
of Section~\ref{sec:variants} say how often that situation arises: \texttt{POLE}
carries $10\,737$ records of uncertain significance against $1891$
pathogenic-matching ones.

\textbf{First step.} A retrospective cohort stratified by MMR/MSI status, with
exonuclease-domain \texttt{POLE} genotype as the exposure and TMB as the
outcome. The analysis must be pre-registered, because the interaction term is
exactly the kind of result that appears in any large cohort if enough subgroups
are tried.

\textbf{What would refute the hypothesis.} No interaction term distinguishable
from zero at a cohort size adequate to detect the effect that
Section~\ref{sec:mmr}'s channel model predicts. Stating that threshold in
advance is part of the first step, not an afterthought.

\subsection*{D6. Formalise what the refutation does \emph{not} reach}

\textbf{The question.} Can the class of models excluded by PSR$_2$ be
characterised exactly, rather than exhibited?

\textbf{Where it comes from.} Section~\ref{sec:nogo} and Exercise~3.3. The
no-go theorem bounds a family; it does not describe the boundary of that family.

\textbf{First step.} Formalise the notion of empirical content used informally
in Exercise~3.3: for a family $\{M_\theta\}_{\theta \in \Theta}$, the set of
observations inconsistent with every member. Then prove the monotonicity ---
enlarging $\Theta$ shrinks that set --- and apply it to the chain
independent~$\subset$~interacting~$\subset$~arbitrary. This is a modest
formalisation with an immediate use: it would let the book state precisely how
much empirical content was given up when $\Lambda$ was introduced in
\S\ref{sub:interact}, rather than observing qualitatively that some was.

\textbf{What would make it pointless.} If the characterisation turns out to be
trivial --- that the excluded set is simply $\{\rho > 1\}$ and nothing more is
provable --- then the exercise is a definition, not a theorem, and should be
written as one.

\subsection*{D7. Audit the rest of this repository the way Section~11 audits the book}

\textbf{The question.} How many of the claims elsewhere in this project are
forced by construction, unsourced, or untested against a null?

\textbf{Where it comes from.} Appendix~\ref{app:repro}, and the pattern of this
edition. The internal review found four errors of fact and three overclaims in
eleven sections. There is no reason to think the rate is lower in material that
has not been reviewed, and some reason to think it is higher.

\textbf{First step.} The three checks that actually found things here, applied
mechanically: \emph{(i)} for every quoted number, locate the primary source and
recompute; \emph{(ii)} for every control, construct the observation that would
have failed it, and if none exists relabel it a consistency check;
\emph{(iii)} for every author--year citation, verify the publication exists.
Check \emph{(iii)} has already found one fabricated reference in this project,
and check \emph{(ii)} has found several forced results.

\textbf{What would make this unnecessary.} Nothing, and that is the point. This
is maintenance rather than research, it produces no new findings, and it is the
highest-value item on the list by any measure that counts avoided error. We
record it last because it is the least interesting and the most likely to be
skipped.

\subsection*{A closing note on how to choose}

The directions above are not equally good science, and the ranking encodes a
judgement worth making explicit. D1 and D4 are measurements that would settle
something; D2 and D6 are theory that would sharpen something; D3 is a method
question whose most likely answer is negative; D5 is the only one with a patient
at the end of it; D7 is housekeeping that will find real errors.

If forced to pick one, pick D1. It is a single measurement, the strains exist,
the technique is established, and the result discriminates between three live
hypotheses including the one this book argues for. The worst outcome is that it
proves us right.
